% GENERATED FILE. Edit canonical lessons, then run npm run generate:books.
% canonical-source-hash: fnv1a64:99993e24d453d9bd
% canonical-lessons: KA-C373-udaya, KA-C373-patana, KA-C373-gatakala, KA-C373-vrttanta, KA-C373-vijaya

\chapter{The Rise and Fall of a Kingdom}
\label{ch:ka-the-rise-and-fall-of-a-kingdom}
\hlchaptermodality{373}

\begin{chapteropening}
\textbf{By the end of this chapter:} \emph{I can say a rise, a downfall, the past, a narrative and a triumph.}

Complete the last lesson of chapter 373: I can say a rise, a downfall, the past, a narrative and a triumph.
\end{chapteropening}

\section[\texorpdfstring{\kn{ಉದಯ} (udaya)}{ಉದಯ (udaya)}]{\kn{ಉದಯ} (udaya) — a rise; a dawn}

\label{lesson:KA-C373-udaya}



\begin{quote}
Before the new one: say the Kannada for former, then the Kannada for theft.
\end{quote}

\subsection*{You'll want to know: \kn{ಉದಯ}}
\textbf{\kn{ಉದಯ}} — \emph{udaya} — \textquotedblleft{}a rise; a dawn\textquotedblright{}.

\begin{quote}
\kn{ಆ} \kn{ರಾಜ್ಯದ} \kn{ಉದಯ} \kn{ಒಂದು} \kn{ಸಣ್ಣ} \kn{ಹಳ್ಳಿಯಲ್ಲಿ} \kn{ಆಯಿತು}.
\end{quote}

\begin{quote}
\emph{ā rājyada udaya ondu saṇṇa haḷḷiyalli āyitu.}
\end{quote}

\begin{quote}
The rise of that kingdom began in a small village.
\end{quote}

\kn{ಸೂರ್ಯೋದಯ} is sunrise. Of a kingdom or a leader, its rise.

\begin{cousinweb}
\textbf{\kn{ಉದಯ}} (\emph{udaya}) is borrowed from Sanskrit \emph{udaya}, \textquotedblleft{}rise, rising of the sun\textquotedblright{} (CDIAL 1931).
\end{cousinweb}

\begin{grammarlens}[title={What you've built}]
One more word for telling what happened.
\end{grammarlens}

\subsection*{Your turn}
\begin{itemize}
\raggedright
  \item \emph{Say it:} \emph{udaya}
  \item \emph{Say it:} \emph{udaya}, once more
  \item \emph{Say it:} \emph{ā rājyada udaya ondu saṇṇa haḷḷiyalli āyitu.} — The rise of that kingdom began in a small village.
  \item \emph{Recall:} say the Kannada for former, then the Kannada for theft, then say \emph{udaya} again
\end{itemize}

\begin{tcolorbox}[breakable,colback=teal!4,colframe=teal!35!black,title={Before you move on}]
What does \kn{ಉದಯ} mean? (\textquotedblleft{}A rise; a dawn\textquotedblright{}.) Say it once more.
\end{tcolorbox}
\section[\texorpdfstring{\kn{ಪತನ} (patana)}{ಪತನ (patana)}]{\kn{ಪತನ} (patana) — a downfall; a fall}

\label{lesson:KA-C373-patana}



\begin{quote}
Before the new one: say the Kannada for theft, then the Kannada for a rise.
\end{quote}

\subsection*{You'll want to know: \kn{ಪತನ}}
\textbf{\kn{ಪತನ}} — \emph{patana} — \textquotedblleft{}a downfall; a fall\textquotedblright{}.

\begin{quote}
\kn{ಸೋಮಾರಿ} \kn{ರಾಜರಿಂದ} \kn{ಆ} \kn{ಸಾಮ್ರಾಜ್ಯದ} \kn{ಪತನ} \kn{ಆಯಿತು}.
\end{quote}

\begin{quote}
\emph{sōmāri rājarinda ā sāmrājyada patana āyitu.}
\end{quote}

\begin{quote}
That empire fell because of lazy kings.
\end{quote}

The opposite of \kn{ಉದಯ}: the fall of a kingdom, a government or a price.

\begin{cousinweb}
\textbf{\kn{ಪತನ}} (\emph{patana}) is borrowed from Sanskrit \emph{patana}, \textquotedblleft{}the act of flying or coming down, falling\textquotedblright{} (CDIAL 7724).
\end{cousinweb}

\begin{grammarlens}[title={What you've built}]
One more word for telling what happened.
\end{grammarlens}

\subsection*{Your turn}
\begin{itemize}
\raggedright
  \item \emph{Say it:} \emph{patana}
  \item \emph{Say it:} \emph{patana}, once more
  \item \emph{Say it:} \emph{sōmāri rājarinda ā sāmrājyada patana āyitu.} — That empire fell because of lazy kings.
  \item \emph{Recall:} say the Kannada for theft, then the Kannada for a rise, then say \emph{patana} again
\end{itemize}

\begin{tcolorbox}[breakable,colback=teal!4,colframe=teal!35!black,title={Before you move on}]
What does \kn{ಪತನ} mean? (\textquotedblleft{}A downfall; a fall\textquotedblright{}.) Say it once more.
\end{tcolorbox}
\section[\texorpdfstring{\kn{ಗತಕಾಲ} (gatakāla)}{ಗತಕಾಲ (gatakāla)}]{\kn{ಗತಕಾಲ} (gatakāla) — the past; bygone days}

\label{lesson:KA-C373-gatakala}



\begin{quote}
Before the new one: say the Kannada for a rise, then the Kannada for a downfall.
\end{quote}

\subsection*{You'll want to know: \kn{ಗತಕಾಲ}}
\textbf{\kn{ಗತಕಾಲ}} — \emph{gatakāla} — \textquotedblleft{}the past; bygone days\textquotedblright{}.

\begin{quote}
\kn{ಅಜ್ಜ} \kn{ಯಾವಾಗಲೂ} \kn{ಗತಕಾಲದ} \kn{ವೈಭವದ} \kn{ಬಗ್ಗೆ} \kn{ಮಾತಾಡುತ್ತಾರೆ}.
\end{quote}

\begin{quote}
\emph{ajja yāvāgalū gatakālada vaibhavada bagge mātāḍuttāre.}
\end{quote}

\begin{quote}
Grandfather always talks about the glory of the past.
\end{quote}

\kn{ಗತ}, gone, and \kn{ಕಾಲ}, time.

\begin{cousinweb}
\textbf{\kn{ಗತಕಾಲ}} (\emph{gatakāla}) is borrowed from Sanskrit, a compound of \emph{gata}, \textquotedblleft{}gone, past\textquotedblright{} (CDIAL 4008), and \emph{kāla}, \textquotedblleft{}time\textquotedblright{} (Monier-Williams).
\end{cousinweb}

\begin{grammarlens}[title={What you've built}]
One more word for telling what happened.
\end{grammarlens}

\subsection*{Your turn}
\begin{itemize}
\raggedright
  \item \emph{Say it:} \emph{gatakāla}
  \item \emph{Say it:} \emph{gatakāla}, once more
  \item \emph{Say it:} \emph{ajja yāvāgalū gatakālada vaibhavada bagge mātāḍuttāre.} — Grandfather always talks about the glory of the past.
  \item \emph{Recall:} say the Kannada for a rise, then the Kannada for a downfall, then say \emph{gatakāla} again
\end{itemize}

\begin{tcolorbox}[breakable,colback=teal!4,colframe=teal!35!black,title={Before you move on}]
What does \kn{ಗತಕಾಲ} mean? (\textquotedblleft{}The past; bygone days\textquotedblright{}.) Say it once more.
\end{tcolorbox}
\section[\texorpdfstring{\kn{ವೃತ್ತಾಂತ} (vṛttānta)}{ವೃತ್ತಾಂತ (vṛttānta)}]{\kn{ವೃತ್ತಾಂತ} (vṛttānta) — a narrative; an account}

\label{lesson:KA-C373-vrttanta}



\begin{quote}
Before the new one: say the Kannada for a downfall, then the Kannada for the past.
\end{quote}

\subsection*{You'll want to know: \kn{ವೃತ್ತಾಂತ}}
\textbf{\kn{ವೃತ್ತಾಂತ}} — \emph{vṛttānta} — \textquotedblleft{}a narrative; an account\textquotedblright{}.

\begin{quote}
\kn{ಪ್ರವಾಸಿ} \kn{ತನ್ನ} \kn{ಪ್ರಯಾಣದ} \kn{ವೃತ್ತಾಂತ} \kn{ಬರೆದನು}.
\end{quote}

\begin{quote}
\emph{pravāsi tanna prayāṇada vṛttānta baredanu.}
\end{quote}

\begin{quote}
The traveller wrote an account of his journey.
\end{quote}

A full telling of what happened, in order.

\begin{cousinweb}
\textbf{\kn{ವೃತ್ತಾಂತ}} (\emph{vṛttānta}) is borrowed from Sanskrit \emph{vṛttānta}, \textquotedblleft{}the end of a course of action, an occurrence, a tale\textquotedblright{} (Monier-Williams).
\end{cousinweb}

\begin{grammarlens}[title={What you've built}]
One more word for telling what happened.
\end{grammarlens}

\subsection*{Your turn}
\begin{itemize}
\raggedright
  \item \emph{Say it:} \emph{vṛttānta}
  \item \emph{Say it:} \emph{vṛttānta}, once more
  \item \emph{Say it:} \emph{pravāsi tanna prayāṇada vṛttānta baredanu.} — The traveller wrote an account of his journey.
  \item \emph{Recall:} say the Kannada for a downfall, then the Kannada for the past, then say \emph{vṛttānta} again
\end{itemize}

\begin{tcolorbox}[breakable,colback=teal!4,colframe=teal!35!black,title={Before you move on}]
What does \kn{ವೃತ್ತಾಂತ} mean? (\textquotedblleft{}A narrative; an account\textquotedblright{}.) Say it once more.
\end{tcolorbox}
\section[\texorpdfstring{\kn{ವಿಜಯ} (vijaya)}{ವಿಜಯ (vijaya)}]{\kn{ವಿಜಯ} (vijaya) — a triumph; a victory}

\label{lesson:KA-C373-vijaya}



\begin{quote}
Before the new one: say the Kannada for the past, then the Kannada for a narrative.
\end{quote}

\subsection*{You'll want to know: \kn{ವಿಜಯ}}
\textbf{\kn{ವಿಜಯ}} — \emph{vijaya} — \textquotedblleft{}a triumph; a victory\textquotedblright{}.

\begin{quote}
\kn{ಯುದ್ಧದಲ್ಲಿ} \kn{ವಿಜಯದ} \kn{ನಂತರ} \kn{ರಾಜ} \kn{ಊರಿಗೆ} \kn{ಹಿಂತಿರುಗಿದನು}.
\end{quote}

\begin{quote}
\emph{yuddhadalli vijayada nantara rāja ūrige hintirugidanu.}
\end{quote}

\begin{quote}
After the victory in the war the king returned to the town.
\end{quote}

\kn{ವಿಜಯ} is the grand word; \kn{ಗೆಲುವು} is the everyday one.

\begin{cousinweb}
\textbf{\kn{ವಿಜಯ}} (\emph{vijaya}) is borrowed from Sanskrit \emph{vijaya}, \textquotedblleft{}contest for victory, victory, conquest, triumph\textquotedblright{} (Monier-Williams), from \emph{ji}, \textquotedblleft{}to win\textquotedblright{}.
\end{cousinweb}

\begin{grammarlens}[title={What you've built}]
One more word for telling what happened.
\end{grammarlens}

\subsection*{Your turn}
\begin{itemize}
\raggedright
  \item \emph{Say it:} \emph{vijaya}
  \item \emph{Say it:} \emph{vijaya}, once more
  \item \emph{Say it:} \emph{yuddhadalli vijayada nantara rāja ūrige hintirugidanu.} — After the victory in the war the king returned to the town.
  \item \emph{Recall:} say the Kannada for the past, then the Kannada for a narrative, then say \emph{vijaya} again
\end{itemize}

\begin{tcolorbox}[breakable,colback=teal!4,colframe=teal!35!black,title={Before you move on}]
What does \kn{ವಿಜಯ} mean? (\textquotedblleft{}A triumph; a victory\textquotedblright{}.) Say it once more.
\end{tcolorbox}
